\begin{figure}[H]
\centering
\begin{subfigure}[t]{0.45\textwidth}
\centering
\begin{tikzpicture}[scale=2.2]
  \draw[->,gray] (-0.3,0) -- (1.35,0) node[right] {$\hat x$};
  \draw[->,gray] (0,-0.9) -- (0,1.2) node[above] {$\hat y$};
  \draw[dashed,gray!60] (0,0) circle (1);
  \draw[-{Stealth},very thick,uniblue] (0,0) -- (1,0);
  \draw[-{Stealth},very thick,unired] (0,0) -- ({cos(40)},{-sin(40)})
        node[below right] {$\hat x'$};
  \draw[unigreen,thick] (0.45,0) arc (0:-40:0.45);
  \node[unigreen] at (0.62,-0.13) {$\theta$};
  \draw[dotted,unired] ({cos(40)},{-sin(40)}) -- ({cos(40)},0);
  \draw[dotted,unired] ({cos(40)},{-sin(40)}) -- (0,{-sin(40)});
  \node[font=\scriptsize,unired,above] at ({cos(40)},0) {$\cos\theta$};
  \node[font=\scriptsize,unired,left] at (0,{-sin(40)}) {$-\sin\theta$};
\end{tikzpicture}
\caption{$\hat x'=\hat x\cos\theta-\hat y\sin\theta$ (componentes sobre $\hat x,\hat y$).}
\end{subfigure}\hfill
\begin{subfigure}[t]{0.52\textwidth}
\centering
\begin{tikzpicture}
\begin{axis}[width=\linewidth,height=6.2cm,xlabel={$\theta$},
  xmin=0,xmax=7,ymin=-2.2,ymax=2.2,
  legend style={font=\scriptsize,at={(0.5,-0.28)},anchor=north,legend columns=3},
  grid=both,minor tick num=1,
  xtick={0,1.5708,3.1416,4.7124,6.2832},
  xticklabels={$0$,$\pi/2$,$\pi$,$3\pi/2$,$2\pi$}]
  \addplot[black,very thick,domain=0:7,samples=200] {cos(deg(x))};
  \addlegendentry{$\cos\theta$}
  \addplot[uniblue,dashed,thick,domain=0:7,samples=100] {1};
  \addlegendentry{orden 0}
  \addplot[unired,thick,domain=0:7,samples=100] {1-x^2/2};
  \addlegendentry{hasta orden 2}
  \addplot[unigreen,thick,domain=0:7,samples=100] {1-x^2/2+x^4/24};
  \addlegendentry{hasta orden 4}
  \addplot[orange,thick,domain=0:7,samples=100] {1-x^2/2+x^4/24-x^6/720};
  \addlegendentry{hasta orden 6}
\end{axis}
\end{tikzpicture}
\caption{Sumas parciales de la serie de Hadamard (coeficiente de $\hat x$).}
\end{subfigure}
\caption{Rotaci\'on finita generada por $\hat L_z$: cada orden de la serie de Hadamard a\~nade un t\'ermino de la serie de Taylor de $\cos\theta$ y $\sin\theta$.}
\label{fig:rot_hadamard}
\end{figure}